\begin{ex}
(UFRGS) Para colocar preço em seus produtos, uma empresa desenvolveu um sistema simplificado de código de barras formado por cinco colunas separadas por quatro espaços. Podem ser usadas colunas de três larguras possíveis e espaços de duas larguras possíveis, como no exemplo abaixo:
\begin{center}
    \begin{tikzpicture}
     \draw [very thick] (0,0)--(0,4); \draw [fill] (.2,0) rectangle (.4,4); \draw [fill] (.8,0) rectangle (1.2,4); \draw [very thick] (1.6,0)--(1.6,4);\draw [fill] (1.8,0) rectangle (2,4);
    \end{tikzpicture}
\end{center}
O número total de preços que podem ser representados por esse código é:
   \begin{enumerate}[(a)]
   \item 1440
   \item 2880
   \item 3125
   \item 3888
   \item 4320
   \end{enumerate}
     \begin{sol}
        resposta: d \\
        $3^5\cdot2^4=3888$
     \end{sol}
\end{ex}